\documentclass[a4paper]{article}

% Español (México) con XeLaTeX: fontspec para fuentes Unicode, polyglossia
% para la división silábica y los nombres del documento en la variante mexicana.
\usepackage{fontspec}
\usepackage{polyglossia}
\setdefaultlanguage[variant=mexican]{spanish}
% Los nombres chinos van como «pinyin (original)»: xeCJK da a los caracteres
% Han su propia fuente; sin ella XeLaTeX los omite con un aviso.
% FandolSong no cubre algunos caracteres de nombres (U+73FA); WenQuanYi Zen Hei sí.
\usepackage[AutoFallBack=true]{xeCJK}
\setCJKmainfont{FandolSong-Regular.otf}
\setCJKfallbackfamilyfont{\CJKrmdefault}{WenQuanYi Zen Hei}
% polyglossia no traduce los nombres de listings.
\AtBeginDocument{\providecommand{\lstlistingname}{}\renewcommand{\lstlistingname}{Listado}%
  \providecommand{\lstlistlistingname}{}\renewcommand{\lstlistlistingname}{Índice de listados}}
\usepackage{amsmath,amssymb}
\usepackage{graphicx}
\usepackage[margin=2.5cm]{geometry}
\usepackage[most]{tcolorbox}
\usepackage{etoolbox}
\usepackage{subcaption}
\usepackage{listings}
\usepackage{hyperref}
\usepackage{booktabs}
\usepackage{float}

\newtcolorbox{knowledgebox}[1]{
    enhanced, colback=blue!5!white, colframe=blue!75!black, colbacktitle=blue!75!black,
    coltitle=white, fonttitle=\bfseries, title={#1}, boxrule=1pt, sharp corners
}

\newtcolorbox{importantbox}[1]{
    enhanced, colback=yellow!10!white, colframe=yellow!80!black, colbacktitle=yellow!80!black,
    coltitle=black, fonttitle=\bfseries, title={#1}, sharp corners
}

\newtcolorbox{warningbox}[1]{
    enhanced, colback=red!5!white, colframe=red!75!black, colbacktitle=red!75!black,
    coltitle=white, fonttitle=\bfseries, title={#1}, sharp corners
}

\lstset{
    basicstyle=\ttfamily\small,
    keywordstyle=\color{blue},
    stringstyle=\color{red!60!black},
    commentstyle=\color{green!60!black},
    breaklines=true,
    frame=single,
    numbers=left,
    numberstyle=\tiny\color{gray},
    captionpos=b,
    extendedchars=true
}

\newcommand{\notetitle}{CS224R Lecture 2: Aprendizaje por imitación}
\newcommand{\noteauthors}{Elaborado a partir de la clase de Chelsea Finn y del material de Stanford CS224R}
\newcommand{\notedate}{9 de abril de 2,025}
\newcommand{\videochannel}{Stanford Online}
\newcommand{\videopublishdate}{2025-04}
\newcommand{\videoduration}{alrededor de 80 minutos}
\newcommand{\videourl}{https://cs224r.stanford.edu/spring\_2025/}
\newcommand{\videocoverpath}{cover.jpg}
\newcommand{\repourl}{https://github.com/hqhq1025/ai-course-notes}


% Footer with repo info
\usepackage{fancyhdr}
\pagestyle{fancy}
\fancyhf{}
\fancyfoot[C]{\thepage}
\fancyfoot[R]{\footnotesize\color{gray}\href{\repourl}{AI Course Notes}}
\renewcommand{\headrulewidth}{0pt}
\renewcommand{\footrulewidth}{0pt}

\begin{document}

\begin{titlepage}
\centering
{\Large Notas del curso\par}
\vspace{1.2cm}
{\huge\bfseries \notetitle\par}
\vspace{0.8cm}
{\large \noteauthors\par}
\vspace{0.3cm}
{\large \notedate\par}
\vspace{1.1cm}

\ifdefempty{\videocoverpath}{
{\small Imagen de portada no proporcionada.\par}
}{
\includegraphics[width=0.82\textwidth,height=0.45\textheight,keepaspectratio]{\videocoverpath}\par
}

\vfill
\begin{tcolorbox}[width=0.9\textwidth, colback=black!2!white, colframe=black!60, sharp corners]
\textbf{Autor o canal del video}: \videochannel\par
\textbf{Fecha de publicación}: \videopublishdate\par
\textbf{Duración del video}: \videoduration\par
\textbf{Ponente}: Chelsea Finn\par
\textbf{Página del curso}: \href{https://cs224r.stanford.edu/spring_2025/}{cs224r.stanford.edu}\par
\textbf{Página del proyecto}: \href{\repourl}{\nolinkurl{\repourl}}\par
\end{tcolorbox}
\end{titlepage}

\tableofcontents
\newpage

\section{Ideas básicas del aprendizaje por imitación}

El aprendizaje por imitación (Imitation Learning) parte de copiar directamente al experto, en lugar de diseñar una función de recompensa compleja. Es adecuado para escenarios con datos suficientes y demostraciones de alta calidad del experto, pero donde el reward no está bien definido.

\begin{importantbox}{El objetivo del aprendizaje por imitación}
Dado un conjunto de trayectorias de demostración $\mathcal{D} = \{(s_1,a_1,\ldots,s_T)\}$ generado por una política experta $\pi_{\text{expert}}$, el objetivo es encontrar una política parametrizada $\pi_\theta$ cuya distribución de acciones predicha coincida lo más posible con la del experto en los mismos estados, de modo que se conserve un desempeño al nivel del experto.
\end{importantbox}

\subsection{Version 0: Regresión de política determinista}

El enfoque más directo es tratar el aprendizaje por imitación como un mapeo de comportamiento:
$$\min_\theta \frac{1}{|\mathcal{D}|}\sum_{(s,a) \in \mathcal{D}} \|\pi_\theta(s) - a\|^2.$$
Este método depende por completo de que los datos de demostración estén bien balanceados y es muy sensible al noise.

\begin{warningbox}{Problema de la media y conductas irreversibles}
Cuando los datos del experto presentan multimodalidad (por ejemplo, girar a la izquierda o a la derecha), el modelo de regresión produce una acción «promedio», lo cual provoca que el sistema entre en estados peligrosos dentro del entorno; este problema es habitual en tareas como la conducción urbana y la manipulación con brazos dobles.
\end{warningbox}

\subsection{Clonación de comportamiento vs. aprendizaje por refuerzo inverso}

La clonación de comportamiento modela directamente las acciones, mientras que el aprendizaje por refuerzo inverso (Inverse RL) reconstruye el reward y después realiza la optimización mediante RL. A continuación se compara el trade-off entre ambos enfoques:

\begin{table}[H]
\centering
\begin{tabular}{p{0.25\textwidth} p{0.32\textwidth} p{0.32\textwidth}}
\toprule
Método & Ventajas & Desventajas \\
\midrule
Clonación de comportamiento & Entrenamiento sencillo, usa directamente supervised learning & Pierde fácilmente la multimodalidad y se ve muy afectada por el cambio de distribución \\
RL inverso & Puede recuperar el reward latente y cuenta con capacidad de generalización & Requiere una gran cantidad de muestreo; el aprendizaje del reward está acoplado con el entrenamiento de la policy \\
\bottomrule
\end{tabular}
\caption{Comparación entre la clonación de comportamiento y el RL inverso}
\end{table}

\begin{knowledgebox}{Por qué priorizar la clonación de comportamiento}
En muchos escenarios industriales (como la conducción autónoma), las trayectorias del experto son suficientes y de alta calidad; en ese caso, aplicar directamente la clonación de comportamiento reduce en gran medida el ciclo de lanzamiento. Si aparece deriva de la distribución, se introduce gradualmente un modelo de reward.
\end{knowledgebox}

\subsection{Puesta en producción en escenarios reales y caution}

Las estrategias habituales en el despliegue real incluyen: el sobremuestreo de los datos de conducción, entrenar una fallback policy usando a los experts como baseline, y establecer un mecanismo de monitoreo human-in-the-loop. En escenarios de fallas de baja frecuencia, es un diseño habitual usar primero la clonación de comportamiento para hacer un warm-start y después aplicar fine-tuning con RL.

\begin{importantbox}{La interpretabilidad de la clonación de comportamiento}
Debido a que la salida es deterministic/log-prob, resulta fácil hacer trace-back; en cuanto el modelo produce una salida anómala, el equipo de ingeniería puede revisar directamente el expert sample más reciente para ubicar rápidamente el origen del problema.
\end{importantbox}

%-\subsection{Resumen de la sección}
-
%-La clonación de comportamiento ofrece...
%-\subsection{Fusión de múltiples expertos y pesos dinámicos}
-
-En la era de los grandes modelos...

\begin{knowledgebox}{El valor de la mezcla de expertos}
Colocar los logits de múltiples expertos dentro de un marco Mixture-of-Experts permite lograr una transferencia zero-shot con pocas muestras, porque la gating network puede orientarse automáticamente hacia el expert más relevante para la tarea actual.
\end{knowledgebox}

\subsection{Resumen de la sección}

La clonación de comportamiento ofrece la ruta de entrada más rápida, pero la regresión deterministic debe combinarse con gating de múltiples expertos y modelado de distribución para poder generalizar en escenarios complejos. La mezcla de múltiples expertos ofrece una asignación automática entre distintas skill, lo que amplía los límites de aplicabilidad del aprendizaje por imitación.

\section{Expresividad de la política y modelos generativos}

Un elemento clave del aprendizaje por imitación es cómo expresar la distribución de la política. Una regresión gaussiana única solo puede representar un comportamiento local y unimodal, y no puede igualar datos de expertos multimodales.

\subsection{Capacidad de representación de distribuciones}

Las acciones discretas pueden modelarse directamente con una distribución categorical, mientras que en el nivel de acciones continuas las formulaciones habituales también incluyen normalizing flow, Mixture Density Network y diffusion model.

\begin{importantbox}{El salto de expresividad que aportan los modelos generativos}
Al considerar la política como una red generativa $p_\theta(a \mid s)$, es posible retomar arquitecturas como diffusion, flow y VAE, de modo que la policy misma se convierta en un modelo probabilístico muestreable e inferible, lo que la adapta de manera natural a datos de demostración multimodales.
\end{importantbox}

\subsection{Políticas Diffusion, Mixture y Autoregressive}

Los medios de expresión más comunes:
\begin{enumerate}
    \item \textbf{Mixture of Gaussians}: produce varios $(\mu_i, \sigma_i, \alpha_i)$ a la vez y los combina con pesos softmax.
    \item \textbf{Autoregressive discretization}: divide el espacio de acciones en subdimensiones y modela de forma recursiva la distribución de cada dimensión.
    \item \textbf{Diffusion Policy}: primero muestrea ruido y después lo elimina de forma iterativa hasta generar la acción final.
\end{enumerate}

\begin{table}[H]
\centering
\begin{tabular}{p{0.22\textwidth} p{0.34\textwidth} p{0.34\textwidth}}
\toprule
Método & Ventaja & Escenario típico \\
\midrule
MoG & Ofrece un modo de mezcla limitado e interpretable & Control continuo de baja dimensión \\
Autoregressive & Descompone las dimensiones y captura correlation de grano fino & Secuencias de comportamiento de conducción autónoma \\
Diffusion & Permite implicit sampling & Manipulación robótica de alta dimensión \\
\bottomrule
\end{tabular}
\caption{Arquitecturas típicas de políticas expresivas}
\end{table}

\subsection{Muestreo y diagnóstico de cobertura}

Una vez entrenados los parámetros de la política, es necesario verificar su cobertura de la expert distribution. Las prácticas habituales incluyen:
\begin{itemize}
    \item KL divergence sweep: estima la KL con datos de expertos y datos muestreados.
    \item Reject sampling: si durante el rollout la probabilidad de la acción cae por debajo del threshold, se activa un human override.
    \item Latent space interpolation: interpola en las variables de control latent y observa si la acción generada cae dentro del manifold de los expertos.
\end{itemize}

\begin{knowledgebox}{Cambio de estrategia de muestreo}
Archit recuerda: \textit{«la estrategia de muestreo determina si podemos observar comportamientos extremos»}. Antes del despliegue, los ajustes de temperature y top-k/top-p de la policy deben escribirse en el runbook y registrarse en el observational log.
\end{knowledgebox}

Flow-based policy se usa para asignar el espacio de acciones a un latent space más regular; durante el entrenamiento se ajusta directamente la distribución objetivo en el latent, de modo que al muestrear basta con avanzar en la dirección latent y después recuperar la acción mediante un invertible map. Para acciones de alta dimensión (como el control humanoid), flow puede reducir considerablemente la dificultad del muestreo.

\begin{importantbox}{La intuición del control latente}
Flow policy equivale a hacer primero un muestreo gaussiano simple en el espacio latente y después convertirlo en una acción compleja mediante un invertible mapping. Esto permite trasladar el objetivo del aprendizaje por imitación al latent space, lo que simplifica la propagación del gradiente.
\end{importantbox}

\subsection{Resumen de la sección}

La expresividad de la política determina si el aprendizaje por imitación puede cubrir el expert space. Modelar la distribución multimodal con modelos generativos, junto con Flow-based latent control y rejection sampling, permite corregir de manera continua en el extremo del rollout, aproximarse a la expert distribution y ofrecer un punto de partida estable para el fine-tuning de downstream RL.

\section{Acumulación de errores y corrección en línea}

\subsection{Origen de los compounding errors}

El challenge central del aprendizaje por imitación consiste en que la acción predicha $\hat{a}$ modifica el siguiente estado $s'$, lo cual produce una distribución de entrada no i.i.d. Los errores pequeños se amplifican progresivamente, hasta producir una catastrophic failure.

\begin{warningbox}{Desplazamiento de distribución y covariate shift}
Que la distribución de estados del experto $p_{\text{expert}}(s)$ y la distribución de la política aprendida $p_{\pi}(s)$ no coincidan es un riesgo sistemático: si no se interviene, aun cuando la loss promedio de entrenamiento sea baja, el performance puede desplomarse en unos cuantos estados.
\end{warningbox}

\subsection{DAgger y Human-Gated DAgger}

La idea de solución consiste en dejar que la política se ejecute por sí misma, recolectar los nuevos estados y usar al experto para proporcionar acciones de corrección en esos estados:
\begin{enumerate}
    \item Ejecutar un rollout de la política actual $\pi_\theta$ y registrar los estados visitados.
    \item Solicitar la acción del experto en cada estado y agregarla al dataset.
    \item Reentrenar la política o actualizar los parámetros mediante importance weighting.
\end{enumerate}

\begin{table}[H]
\centering
\begin{tabular}{p{0.28\textwidth} p{0.28\textwidth} p{0.28\textwidth}}
\toprule
Etapa & Entregable clave & Elemento de revisión \\
\midrule
Observe & Rollout log & state distribution shift \\
Assess & Expert correction & human agreement rate \\
Act & Updated policy checkpoint & performance vs baseline \\
\bottomrule
\end{tabular}
\caption{Entregables y revisiones en el proceso de DAgger}
\end{table}

\begin{knowledgebox}{Offline vs. Online del aprendizaje por imitación}
\begin{itemize}
    \item \textbf{Offline (clonación de comportamiento)}: solo usa un static dataset; los datos son seguros, pero la generalización es deficiente.
    \item \textbf{Online (DAgger/HG-DAgger)}: implica human-in-the-loop; es eficiente, pero requiere permisos de experto.
\end{itemize}
\end{knowledgebox}

\subsection{Sinergia entre el aprendizaje por imitación y el aprendizaje por refuerzo}

En la práctica, es común usar BC como warm-start y luego aplicar RL fine-tune, formando un hybrid pipeline: primero se aprende una policy inicial con demonstrations de alta calidad, y después se mejora la robustez mediante rollout distribuido + reward fine-tuning.

\begin{warningbox}{No abandonar la señal de reward}
Aunque el objetivo sea imitar al experto, también conviene monitorear de manera continua el reward potencial (por ejemplo, safety penalty, human critique score), para poder hacer un rollback de inmediato en cuanto la recompensa presente drift.
\end{warningbox}

\subsection{Resumen de la sección}

Los compounding errors requieren combinar la recolección de datos con estrategias de human-in-the-loop. DAgger/HG-DAgger ofrece una ruta de corrección, y al combinarla con un RL fine-tuning posterior, la policy puede seguir convergiendo bajo distribution shift.

\section{Recopilación y aseguramiento de calidad de los datos de demostración}

\subsection{Pipeline de recopilación de datos}

Los datos de demostración provienen de humanos, robots simulados o pre-recorded log. Pipelines comunes:
\begin{itemize}
    \item \textbf{Kinesthetic Teaching}: el cuerpo humano controla directamente el brazo robótico, con baja latency.
    \item \textbf{Remote teleoperation}: control en cuadrícula, que permite expertos distribuidos geográficamente.
    \item \textbf{Behavioral cloning from screen recording}: sincronización de video con la actividad de teclado y mouse.
\end{itemize}

\begin{knowledgebox}{embodiment gap}
Usar directamente video humano en lugar de robot demonstration produce un embodiment gap (diferencia de grados de libertad, diferencia de dinámica), pero puede servir como exploration signal, y el gap se puede reducir mediante domain randomization.
\end{knowledgebox}

\subsection{QA de datos y audit}

Las demostraciones de alta calidad requieren un audit pipeline, que incluye:
\begin{itemize}
    \item agreement rate entre expertos (Cohen's κ);
    \item variance coverage de los pares estado-acción;
    \item labeling de las demostraciones erróneas;
    \item Metadata (dificultad de la tarea, lighting del entorno).
\end{itemize}

\begin{table}[H]
\centering
\begin{tabular}{p{0.28\textwidth} p{0.35\textwidth} p{0.33\textwidth}}
\toprule
Dimensión & Métrica & Respuesta de ingeniería \\
\midrule
Consistencia & Pairwise agreement & Los ejemplos inconsistentes se envían a revisión \\
Diversidad & Different skills, contexts & Se usa clustering para identificar gaps de cobertura \\
Seguridad & Adversarial prompt & Se agregan counterfactual data \\
\bottomrule
\end{tabular}
\caption{Aseguramiento de calidad de los datos de demostración}
\end{table}

\subsection{Colaboración entre modalidades y entre roles}

Los proyectos modernos de aprendizaje por imitación suelen involucrar a equipos de etiquetado de datos, expertos en lenguaje y SRE, que mantienen en conjunto el dataset catalog y el manifest. Cada rollout requiere agregar new prompts al dataset y registrar el shift log en el evidence board.

\subsection{Resumen de la sección}

La calidad de los datos de demostración determina directamente el límite inferior del aprendizaje por imitación. Mediante un audit riguroso, metadata y multi-role feedback, es posible convertir problemas como el embodiment gap o el adversarial prompt en riesgos gobernables.

\section{Despliegue y confianza}

\subsection{Pila de evaluación y métricas de monitoreo}

Después del despliegue es necesario conectar el evidence de la etapa de entrenamiento con las métricas de producción:
\begin{itemize}
    \item Benchmark (HumanEval, LongBench) garantiza la capacidad;
    \item Guardrail (Hallucination, Forbidden topics) garantiza la seguridad;
    \item Ops metrics (Rejection rate, Latency P99) garantiza la experiencia.
\end{itemize}

\begin{table}[H]
\centering
\begin{tabular}{p{0.25\textwidth} p{0.55\textwidth}}
\toprule
Categoría & Métrica típica \\
\midrule
Benchmark & HumanEval, TruthfulQA, Domain-specific tests \\
Safety & Red team rejection, prompt-based filter precision \\
Ops & Rejection rate, Latency P99, Human override count \\
\bottomrule
\end{tabular}
\caption{Pila de monitoreo tras el despliegue del aprendizaje por imitación}
\end{table}

\subsection{Diapositivas y evidencia visual}

\begin{figure}[H]
\centering
\includegraphics[width=0.8\textwidth]{slide-02-05.jpg}
\caption{Slide 02-05 muestra el ciclo cerrado que conecta los datos de demostración con el monitoreo (data → policy → ops board).}
\end{figure}

\begin{figure}[H]
\centering
\includegraphics[width=0.8\textwidth]{slide-02-10.jpg}
\caption{Slide 02-10 combina el evaluation stack y el governance checkpoint en una sola matriz, y enfatiza el evidence-driven rollout.}
\end{figure}

\begin{importantbox}{El papel de los tableros de visualización}
Visualizar el evidence stacking como un mapa de calor en las diapositivas permite que PM, el equipo legal y SRE discutan los riesgos desde la misma perspectiva de datos, en lugar de depender de documentos dispersos.
\end{importantbox}

\subsection{Retroalimentación y gobernanza}

Se sigue el modelo Observe-Assess-Act-Document:
\begin{itemize}
    \item Observe: se registra el desempeño del rollout y el reward drift;
    \item Assess: se revisa el guardrail con la evidence matrix;
    \item Act: el human-in-the-loop activa el mitigation;
    \item Document: la decisión se registra en el runbook, para una auditoría posterior.
\end{itemize}

\begin{warningbox}{Documentar es más importante que corregir}
Un loop de retroalimentación sin runbook hace que al equipo le resulte difícil recordar por qué se hizo rollback en cierto prompt. Archit sugiere: cada mitigation debe ir acompañado de un drift snapshot, un human override log y un lesson learned.
\end{warningbox}

\subsection{Resumen de la sección}

La confianza en el despliegue proviene del ciclo cerrado entre data, policy y ops. Solo al convertir el evidence dashboard de las diapositivas en tablas reales de monitoreo y gobernanza se puede lograr que el sistema de aprendizaje por imitación se mantenga estable de manera continua en production.

\section{Estudio de caso y diseño experimental}

\subsection{Ejemplo de robótica: pipeline de pick-and-place}

\begin{figure}[H]
\centering
\includegraphics[width=0.8\textwidth]{slide-02-03.jpg}
\caption{La diapositiva 02-03 muestra un pipeline de pick-and-place, desde los datos de demostración y el modelo generativo hasta el rollout monitoring.}
\end{figure}

Dentro de este pipeline se agregan las expert trajectory, se entrena la diffusion policy, y en el scoring server se calcula en tiempo real el KL drift. El post-mortem anterior al despliegue incluye: 1) control length 2) intervention log 3) visual replay diff.

\begin{table}[H]
\centering
\begin{tabular}{p{0.25\textwidth} p{0.7\textwidth}}
\toprule
Etapa & Medición central \\
\midrule
Datos & Coverage of pick/place, human override rate, sensor noise \\
Entrenamiento & Loss plateau, latent interpolation, sampling temperature drift \\
Despliegue & KL drift, rejection rate, human override latency \\
\bottomrule
\end{tabular}
\caption{Puntos de monitoreo del pipeline de aprendizaje por imitación de pick-and-place}
\end{table}

\begin{knowledgebox}{El papel de la reproducción de video}
Grabar los policy rollouts en short clips y compararlos (diff) con el expert clip ayuda al PM a identificar rápidamente el origen del drift, y así decidir si se hace un rollback.
\end{knowledgebox}

\subsection{Evaluation Matrix y experimentos de alineación}

\begin{figure}[H]
\centering
\includegraphics[width=0.8\textwidth]{slide-02-07.jpg}
\caption{La diapositiva 02-07 coloca benchmark, safety y governance en una matriz por capas, para facilitar la alineación cross-team.}
\end{figure}

El Evaluation matrix se compone de tres capas: benchmark scores, safety metrics, ops guardrails. Cada versión debe pasar las siguientes tests:
\begin{itemize}
    \item Benchmark: MMLU-like tasks, LongBench multi-hop.
    \item Safety: forbidden content rejection, hallucination detectors.
    \item Ops: latency, human override, drift ticket count.
\end{itemize}

\begin{warningbox}{El riesgo de sobreajustar el benchmark}
Perseguir el benchmark sin criterio hace que la policy ignore los rare prompt; es indispensable usar safety y ops guardrail como condiciones de gating, para evitar un drift «impulsado por métricas».
\end{warningbox}

\subsection{Resumen de la sección}

El estudio de caso convierte los datos de demostración, la política generativa, el evaluation matrix y la gobernanza en artifacts concretos (dashboard, table, clips), y ofrece un marco experimental replicable para las extensiones futuras.

\section{Gobernanza e interpretabilidad}

\subsection{Drift Playbook}

\begin{figure}[H]
\centering
\includegraphics[width=0.8\textwidth]{slide-02-08.jpg}
\caption{La diapositiva 02-08 muestra el drift response workflow, que incluye alert, human review y rollback.}
\end{figure}

El Drift Playbook incluye: Detect (alert rules), Diagnose (cluster drift prompt), Decide (governance board), Document (runbook entry). Cada respuesta debe ir acompañada de snapshot, la curva de KL y el human override log.

\begin{table}[H]
\centering
\begin{tabular}{p{0.22\textwidth} p{0.65\textwidth}}
\toprule
Etapa & artifact principal \\
\midrule
Detect & Reward drift alert + latency spike dashboard \\
Diagnose & Prompt clustering report + human override video \\
Decide & Governance board minutes + rollback checklist \\
Document & Runbook entry + postmortem draft \\
\bottomrule
\end{tabular}
\caption{Entrega de artifacts del Drift Playbook}
\end{table}

\begin{knowledgebox}{Los criterios de las alertas automatizadas}
Se recomienda usar tres líneas de alerta: KL drift, hallucination flag y human override frequency; la governance review solo se activa cuando las tres exceden el umbral al mismo tiempo, para evitar false positives frecuentes.
\end{knowledgebox}

\subsection{Automated Evaluation Matrix}

Automation coverage incluye unit-level tests (safety prompt injection, allowed domains), system-level evaluation (trajectory success rate) y ops-level metrics (human override latency). Estas métricas deben mostrarse por niveles en el dashboard y someterse a un cross-check con benchmark/human eval.

\begin{importantbox}{El significado multinivel de la evaluation matrix}
Separar benchmark, safety y ops en niveles, y luego representar el estado de salud con un color-coded heatmap, permite que compliance, ops y research compartan una misma percepción de los hechos en una sola review.
\end{importantbox}

\subsection{Supervisión humana y reuniones de gobernanza}

Cada versión debe designar un alignment owner, un ops owner y un legal owner. El alignment owner mantiene la evidence matrix; el ops owner monitorea los drift dashboards; el legal owner decide si un high-risk prompt necesita un guardrail adicional.

\begin{warningbox}{No pase por alto la governance meeting}
Si la governance meeting es solo una formalidad, el equipo terminará alineándose con la línea que le resulte más fácil de cumplir cuando los indicadores entren en conflicto (generalmente el benchmark). Es necesario preparar un evidence scorecard antes de la reunión y publicar los action items después.
\end{warningbox}

\subsection{Resumen de la sección}

El ciclo cerrado de gobernanza e interpretabilidad es una accountability practice: un drift playbook de alta calidad, una automated evaluation matrix y una governance meeting con múltiples roles le dan al deployment de aprendizaje por imitación trazabilidad y capacidad de respuesta rápida.
\section{Depuración y visualización}

\subsection{Panel de control a nivel de prompt}

Después del despliegue, la wrap-around instrumentation incluye: prompt signature fingerprints, action entropy, KL drift trace, human override counter. Escribir estas señales en el time-series dashboard permite ubicar rápidamente qué tipo de prompt empieza a activar drift.

\begin{table}[H]
\centering
\begin{tabular}{p{0.25\textwidth} p{0.65\textwidth}}
\toprule
Métrica & Descripción \\
\midrule
Prompt signature & hash + metadata, empleado para volver a reproducir y reproduce \\
Entropy & una caída de action entropy puede indicar policy collapse \\
KL drift & policy actual frente a reference distribution \\
Override count & número de veces que human reescribe el prompt, que mide automation trust \\
\bottomrule
\end{tabular}
\caption{Señales centrales del panel de control a nivel de prompt}
\end{table}

\begin{importantbox}{Las métricas deben vincularse al runbook}
Cada signal debe tener su threshold escrito en el runbook, con un owner asignado. Si se excede el límite, hay que notificar automáticamente a la persona correspondiente (ops, alignment, legal), en lugar de depender de comentarios en slack.
\end{importantbox}

\subsection{Visualización de rollout}

\begin{figure}[H]
\centering
\includegraphics[width=0.8\textwidth]{slide-02-04.jpg}
\caption{Slide 02-04 muestra de forma sincronizada la rollout trajectory, el human override y el drift indicator, formando un panel de visualización unificado.}
\end{figure}

El panel de visualización permite que el engineer determine en menos de un minuto: 1) si ese rollout se desvía del expert manifold; 2) si el human override se concentra en cierto tipo de prompt; 3) si las curvas de KL/entropy bajan de forma sincronizada. Estos gráficos se embed en el internal dashboard, para usarse como gate review.

\begin{knowledgebox}{Valor diagnóstico de la visualización}
Cuando un rollout presenta drift, la visualization permite que el equipo vea el prompt específico, la trayectoria de KL y la posición del override sin tener que revisar los registros, lo cual reduce considerablemente el tiempo de incident response.
\end{knowledgebox}

\subsection{Resumen de la sección}

La depuración y la visualización concretan las señales abstractas de drift: el Prompt dashboard ofrece monitoreo cuantitativo, mientras que la Rollout visualization aporta un qualitative insight sobre los puntos de activación; solo combinando ambos se puede completar el análisis de causa raíz antes del siguiente release.
\section{Trabajo en equipo y evolución continua}

\subsection{Matriz de roles y responsabilidades}

Un sistema exitoso de aprendizaje por imitación necesita la colaboración de cuatro equipos: alignment, ops, data y legal. La siguiente matriz enumera el output y el checkpoint de cada rol:

\begin{table}[H]
\centering
\begin{tabular}{p{0.22\textwidth} p{0.28\textwidth} p{0.28\textwidth} p{0.1\textwidth}}
\toprule
Rol & Entregable & Checkpoint principal & Frecuencia \\
\midrule
Alignment Engineer & Evidence matrix, policy checkpoint & KL drift < θ, preference score & per model release \\
Ops/SRE & Drift dashboard, latency log & alert resolution time & daily \\
Data Engineer & Demo catalog, metadata & coverage report & weekly \\
Legal/Compliance & Guardrail doc, review minutes & high-risk prompt gating & per sprint \\
\bottomrule
\end{tabular}
\caption{Matriz de responsabilidades del equipo de aprendizaje por imitación}
\end{table}

\begin{warningbox}{No difuminar la responsabilidad}
Si un guardrail determinado no es responsabilidad ni de ops ni de compliance, nadie responderá durante un incident. Hay que usar la matriz para definir con claridad al owner e incorporar el proceso al team handbook.
\end{warningbox}

\subsection{Evolución continua y direcciones de investigación}

En la vía de research, el aprendizaje por imitación puede avanzar hacia sim-to-real, multi-agent collaboration y la inferencia automatizada de reward. La diapositiva 02-11 destaca, dentro de los escenarios multi-agent, el cooperative imitation, el competitive imitation y el dynamic gating.

\begin{importantbox}{Lo esencial de los future experiments}
Cada research experiment debe ir acompañado de una evaluation matrix que confirme, al terminar el experimento, si las métricas vuelven al baseline; si el drift de las métricas es severo, hay que posponer el rollout a producción.
\end{importantbox}

\subsection{Resumen de la sección}

El trabajo en equipo y la research planning forman el doble eje de la evolución continua: la matriz de responsabilidades garantiza que en cada release haya alguien revisando los datos, y las direcciones de investigación, mediante structured experiments, hacen que los pipelines de aprendizaje por imitación sean robustos frente a los futuros escenarios multi-agent y sim-to-real.
\section{Sim-to-Real y experimentos de verificación}

\subsection{Calibración de la simulación y verificación en el mundo real}

\begin{figure}[H]
\centering
\includegraphics[width=0.8\textwidth]{slide-02-11.jpg}
\caption{Slide 02-11: cómo se calibra un pipeline de sim-to-real mediante domain randomization y réplicas del mundo real.}
\end{figure}

Para que la policy migre de la simulación al entorno real, las estrategias habituales incluyen domain randomization, dynamics randomization y fidelity tuning. Cada migración debe registrar los parámetros de randomization y las métricas de monitoreo del reality gap.

\begin{table}[H]
\centering
\begin{tabular}{p{0.3\textwidth} p{0.6\textwidth}}
\toprule
Punto de verificación & Descripción \\
\midrule
Domain gap & distribution drift score de las diferencias de dinámica/percepción \\
Simulator fidelity & error L2 entre la fidelity visual/de control y el real world log \\
Transfer success & task success rate del rollout en el mundo real \\
\bottomrule
\end{tabular}
\caption{Matriz de verificación de Sim-to-Real}
\end{table}

\begin{knowledgebox}{Registrar los metadatos de transfer}
Usar metadata para registrar, en cada transfer, el simulator seed, el randomization range y la hardware configuration ayuda a reproducir el éxito o el failure de un salto determinado.
\end{knowledgebox}

\subsection{Cross-environment evaluation}

Distintos entornos (lab, field, cloud) tienen distintos constraints. Es necesario fijar un evaluation order, por ejemplo: primero ejecutar el benchmark en el lab environment, después revisar el sensor noise y el drift en el field environment, y por último recolectar logs en el cloud pipeline.

\begin{warningbox}{No usar el baseline directamente en field rollouts}
Una vez que el field environment sufre drift, resulta muy expensive. Hay que verificar primero la métrica en el lab, luego hacer una incremental expansion hacia el field, y por último recolectar los logs en el cloud dentro de un evidence board.
\end{warningbox}

\subsection{Resumen de la sección}

El pipeline de Sim-to-Real depende de una structured validation matrix y de metadata para poder migrar la policy con seguridad entre distintos entornos; antes de cada field rollout es indispensable verificar el domain gap.
\section{Aprendizaje continuo y colaboración multi-agent}

\subsection{Manifestación multi-agent e imitation}

\begin{figure}[H]
\centering
\includegraphics[width=0.8\textwidth]{slide-02-12.jpg}
\caption{Slide 02-12 muestra el aprendizaje por imitación y los communication loops en un multi-agent scenario.}
\end{figure}

En un escenario multi-agent, cada agent puede imitar a distintos specialist, y comparten policy logits entre sí mediante una gating network. La CSI gating permite elegir, en runtime, el expert más adecuado.

\begin{importantbox}{El poder de multi-agent imitation}
Mediante multi-agent imitation, el sistema puede dividir el trabajo y colaborar en tareas complejas que un modelo de un solo agent no puede cubrir. Cada agent se encarga de un segmento de trajectory y, al final, un aggregator sintetiza una política consistente.
\end{importantbox}

\subsection{Aprendizaje continuo y automejora}

\begin{figure}[H]
\centering
\includegraphics[width=0.8\textwidth]{slide-02-13.jpg}
\caption{Slide 02-13 dibuja el continuous improvement loop: data → model → evaluation → deployment.}
\end{figure}

Dentro del continuous improvement loop, cada deployment debe tomar el human override como new training data, e incorporar un multi-agent critic como auxiliary loss dentro del nightly retrain pipeline.

\begin{knowledgebox}{Self-supervised critic}
Usar un self-supervised critic permite estimar el reward gap cuando falta human comparison, y combinado con human override forma una stratified preference signal.
\end{knowledgebox}

\subsection{Resumen de la sección}

Multi-agent imitation y el continuous improvement loop convierten el aprendizaje por imitación en un sistema que se autocalibra constantemente; basta con agregar traction guardrails en evaluation → retrain → deployment para lograr una mejora continua.

\section{Síntesis y ampliación}

\begin{enumerate}
    \item La clonación de comportamiento es la vía más rápida para comenzar, pero hay que estar atento a la conducta promedio y al covariate shift.
    \item Los modelos generativos dotan a la política de capacidad de expresión multimodal, y es necesario ajustar KL/temperature en la etapa de muestreo.
    \item DAgger/HG-DAgger/las estrategias híbridas de RL son una vía eficaz para resolver los compounding errors.
    \item El audit, los metadata y la colaboración multi-role de los datos de demostración de alta calidad son la base para la puesta en producción.
    \item La etapa de despliegue debe girar en torno al evidence board, vinculando las tres líneas de benchmark, safety y ops a la iteración de monitoreo.
    \item Hay que registrar cada evento de mitigation y rollback, para que el governance audit tenga fundamento.
\end{enumerate}

\subsection{Hoja de ruta de puesta en producción}

\begin{figure}[H]
\centering
\includegraphics[width=0.8\textwidth]{slide-02-14.jpg}
\caption{La diapositiva 02-14 muestra en forma de timeline los deliverables desde los datos hasta el deployment.}
\end{figure}

La hoja de ruta de puesta en producción divide el pipeline de aprendizaje por imitación en tres etapas: 1) preparación de datos (demo catalog + metadata), 2) entrenamiento de la política (diffusion/MoG + flow calibration), 3) recuperación ante desastres (drift playbook + governance review). Cada etapa debe definir guardrail, evaluation metric y artifacts.

\begin{table}[H]
\centering
\begin{tabular}{p{0.25\textwidth} p{0.55\textwidth}}
\toprule
Etapa & Key artifact \\
\midrule
Preparación de datos & Metadata-rich demo catalog + coverage report \\
Entrenamiento de la política & KL drift sweep + sampling temperature record \\
Recuperación de despliegue & Drift playbook + rollback checklist \\
\bottomrule
\end{tabular}
\caption{Hoja de ruta por etapas para la puesta en producción del aprendizaje por imitación}
\end{table}

\begin{warningbox}{No confundas el roadmap con un plan de proyecto}
El roadmap debe actualizarse constantemente: se obtienen el delay y el drift reales del evidence dashboard, y luego los lesson learned se escriben de vuelta en el timeline, para poder seguir avanzando.
\end{warningbox}

\subsection{Tabla de resumen}

\begin{table}[H]
\centering
\begin{tabular}{p{0.2\textwidth} p{0.35\textwidth} p{0.35\textwidth}}
\toprule
Tema & Takeaway central & Acción de ingeniería \\
\midrule
Clonación de comportamiento & Puesta en marcha rápida, pero hay que atender el covariance shift & Ejecutar un KL sweep después del entrenamiento y un sanity check antes del despliegue \\
Expresión de la política & La distribución generativa necesita explorarse por muestreo & Usar diffusion/MoG y registrar el sampling temperature \\
Corrección en línea & DAgger compensa los estados no vistos & Establecer un runbook de Observe-Assess-Act-Document \\
Control de calidad de datos & Multi-role audit & Usar metadata + clustering para identificar coverage gaps \\
Gobernanza del despliegue & Evidence dashboard & Vincular las tres líneas de benchmark/safety/ops \\
\bottomrule
\end{tabular}
\caption{Síntesis accionable de esta clase}
\end{table}

\subsection{Lecturas adicionales}
\begin{itemize}
    \item Chi et al., \textit{Diffusion Policy}: incorpora el diffusion model a las políticas robóticas.
    \item Ross et al., \textit{A Reduction of Imitation Learning and Structured Prediction to No-Regret Online Learning}: el artículo original de DAgger.
    \item Brohan et al., \textit{RT-2}: muestra los efectos del aprendizaje por imitación multimodal.
    \item Rajeswaran et al., \textit{Behavioral Cloning with Expert Interventions}: sistematiza las human interventions.
\end{itemize}

\subsection{Resumen de la sección}

Esta clase recorre desde la clonación de comportamiento hasta las políticas generativas, y desde el DAgger/online pipeline hasta el data QA y la gobernanza del despliegue, conformando un marco práctico end-to-end para el aprendizaje por imitación. Solo cuando todas las etapas están conectadas entre sí, el aprendizaje por imitación puede ejecutarse de manera verdaderamente confiable en el entorno de producción.

\end{document}
